\documentclass[12font]{article}
\title{\textbf{ Quantum Metric and PSG }}
\author{Chen-Chih Wang  }
\date{\textit{Department of physics, National Tsing Hua University, Hsinchu 30013, Taiwan} \\ (Dated: \today)}
\usepackage[margin=2cm]{geometry}
%\usepackage[fleqn]{amsmath}
%\usepackage{CJK,amsmath}
\usepackage{amsfonts,amssymb}
\usepackage{fancyhdr}
%\pagestyle{fancy}
\usepackage{textcomp}
\usepackage{graphicx}
\usepackage{float}
\usepackage{color}
\usepackage{tikz}
\usepackage{multicol}
\usepackage{threeparttable}
\usepackage{amsthm}
\usepackage{mathtools}
\usepackage{hyperref}
\urlstyle{same}
\usepackage{caption}
\usepackage{subcaption}
\usepackage{dsfont}

\newcommand{\bra}[1]{\left\langle #1 \right|}
\newcommand{\ket}[1]{\left| #1 \right\rangle}
\newcommand{\anglelr}[1]{\left\langle #1 \right\rangle}
\newcommand{\braket}[2]{\left\langle #1 \right|\left. #2 \right\rangle}
\newcommand{\dif}[1]{\frac{\mathrm{d}}{\mathrm{d}#1}}
\newcommand{\diff}[2]{\frac{\mathrm{d}#1}{\mathrm{d}#2}}
\newcommand{\gd}{\boldsymbol{\nabla}}
\newcommand{\dive}{\boldsymbol{\nabla}\cdot}
\newcommand{\curl}{\boldsymbol{\nabla}\times}
\newcommand{\lr}[1]{\left(  #1  \right)}
\newcommand{\lrm}[1]{\left[  #1  \right]}
\newcommand{\lrg}[1]{\left\{  #1  \right\}}
\newcommand{\abs}[1]{\left|  #1  \right|}
\newcommand{\de}{\partial}
\newcommand{\pdif}[1]{\frac{\partial}{\partial#1}}
\newcommand{\pdiff}[2]{\frac{\partial#1}{\partial#2}}
\newcommand{\mrm}[1]{\mathrm{#1}}
\newcommand{\bo}[1]{\mathbf{#1}}
\newcommand{\bog}[1]{\boldsymbol{#1}}
\newcommand{\vint}[3]{\int_{#1}#2\cdot\mathrm{d}\mathbf{#3}}
\newcommand{\voint}[3]{\oint_{#1}#2\cdot\mathrm{d}\mathbf{#3}}
\newcommand{\Vint}{\int\mrm{d}^3\bo{r}'}
\newcommand{\rr}{\mathbf{r} - \mathbf{r}'}
\newcommand{\arr}{\left|\mathbf{r} - \mathbf{r}'\right|}

\begin{document}
\maketitle

%\fontsize{12pt}{18pt}\selectfont


\begin{abstract}
  The quantum metric is highly related to the geometry of the Berry connection. In our theory, the entanglement order parameter in the PSG theory can be considered the Berry connection. Here I am trying to build the connection between the quantum metric and the entanglement pattern, such that we can use the geometric point of view to comprehend the entanglement structure of a spin liquid. 
\end{abstract}

\section{Introduction}
The quantum metric has an important application in QHE and FQHE. And we also know that the spin liquid has the similar behaviour to FQH since both of them have fractional excitation. 



\begin{thebibliography}{99}

    \bibitem{}{}

\end{thebibliography}

\end{document}

